\chapter{Model behaviour fundamentals}
\label{ch:model}

\section*{What this chapter is for}

The universal baseline. \reported{These are asked everywhere and they are not
where anybody is differentiated}{field-guide} \dash{} getting one wrong is
expensive and getting them all right earns nothing. Treat this chapter as
hygiene.

\section*{The questions}

\begin{bankq}{What does temperature actually do, and when would you change it?}
\qasked
Almost always asked, and almost always early. It is a warm-up, which means
getting it wrong is expensive and getting it right earns nothing.

\qtested
Whether you can describe a sampling parameter without reaching for
metaphor. \emph{Creativity} is the word to avoid.

\qstaff
The model produces a distribution over next tokens; temperature reshapes it
before sampling. Low values concentrate the mass on the most likely tokens;
high values flatten it. At zero it is effectively taking the most likely token
each time, which makes runs repeatable in practice while not being a guarantee
across providers, versions or hardware.

When to change it is the useful half, and it is a property of the task rather
than of the product. Extraction, classification, anything feeding a parser:
low, because you want the same answer twice. Anything where several answers
are acceptable and variety is the point: higher. \judgement{A single global
temperature for a system doing both is a mistake people make because it is a
configuration value rather than a per-call decision.}

\qsenior
Low for deterministic tasks, high for creative ones \dash{} correct, and framed
as a product setting rather than a per-call one.

\qcontested
Whether temperature zero gives reproducibility. \judgement{Treat it as reducing
variance rather than removing it; anything that depends on exact
reproducibility needs a check rather than a setting.}

\qdrill
Name three calls in one system that should have different temperatures, and
say why for each.
\end{bankq}

\begin{bankq}{How do you get structured output reliably?}
\qasked
Usually arrives as a practical complaint rather than a question: \emph{we get
JSON back and sometimes it does not parse}.

\qtested
Whether you have shipped something that parses model output. There is a
characteristic set of scars.

\qstaff
Use the provider's constrained mode where it exists \dash{} schema-enforced
output, or a tool call whose parameters are the structure you want. It removes
a whole class of failure rather than reducing it.

Where it does not exist, the order is: ask for the structure explicitly, parse
strictly, and on a parse failure retry with the error rather than with the same
prompt. And validate semantically after parsing \dash{} well-formed and correct
are different, and a schema will happily accept a date in the wrong century.

\judgement{The thing to say that marks experience: never repair the output with
string manipulation. Trimming stray text around a payload works until the day
the stray text is inside it, and by then the repair is load-bearing and
undocumented.}

\qsenior
Ask for it in the prompt, parse it, retry on failure. Works, and skips
constrained decoding, which is the actual answer where it is available.

\qcontested
Little.

\qdrill
Describe what your system does when the third retry also fails to parse. Most
do not have an answer, and that is the finding.
\end{bankq}

\begin{bankq}{What happens as the conversation gets long?}
\qasked
Asked as a design question about a chat product, or as a debugging question
about one that got slow and expensive.

\qtested
Whether you know the window is a budget you spend rather than a limit you hit.

\qstaff
Three things, and only the first is widely known. You approach the limit, and
have to do something \dash{} truncate, summarise, or retrieve. You pay for the
whole history on every turn, so cost grows with the square of the conversation
length unless you intervene. And quality changes with position: material in the
middle of a long context is used less reliably than material at either end,
which means \emph{what} you drop matters less than \emph{where} you put what
you keep.

The practical consequence: put the instruction and the most relevant material
where it will be attended to, and decide the truncation policy deliberately
rather than letting it be whatever the library does.

\qsenior
Naming the limit and summarisation as the fix. Correct on the first of the
three.

\qcontested
How strong the position effect is, and whether it holds across current models,
is disputed and the evidence moves. \judgement{Design as though it holds; the
cost of doing so is close to zero.}

\qdrill
For a system you know, say what happens on the turn where the history no longer
fits. If the answer is ``the library handles it'', find out how.
\end{bankq}

\section*{What this chapter does not claim}

These answers describe common behaviour across providers at the time of
writing. Specific behaviours \dash{} constrained decoding support, how
temperature is implemented, position effects \dash{} vary by provider and
version, and this chapter deliberately names none, because a claim about a
provider's behaviour is a claim about a version.
